\documentclass{article}

% Paquetes
\usepackage[utf8]{inputenc} % Para caracteres en español
\usepackage[spanish]{babel} % Para la correcta separación de sílabas y etiquetas en español
\usepackage{graphicx}
\usepackage{amssymb}
\usepackage{rotating}
\usepackage{float}
\usepackage{hyperref}
\usepackage{xcolor}

\title{}
\author{}
\date{}
\setlength{\parskip}{10pt}


\begin{document}
\maketitle

\section{Estudio de la distribución del tiempo y gestión del hogar en entornos compartidos}

Un estudio del \textit{American Time Use Survey (ATUS)}, analiza la distribución del tiempo que las personas dedican a la administración del hogar
y cómo esta tarea se realiza en intervalos dispersos. Este estudio revela patrones en las actividades domésticas y ayuda a entender los 
desafío de coordinación y planificación en hogares donde las responsabilidades se dividen entre varias personas. El estudio muestra
que las tareas del hogar, como la limpieza, la preparación de comidas, la organización y la gestión de finanzas, tienden a realizarse en 
intervalos de tiempo que varían significativamente entre los distintos miembros del hogar. Esto es problemático en hogares compartidos, como 
los estudiantes o trabajadores, donde la disponibilidad de cada persona depende de su carga académica o laboral. La falta de tiempos fijos y 
la irregularidad en la dedicación a las tareas puede generar problemas de coordinación, ya que cada miembro tiene su propio horario, lo que impide
una rutina doméstica unificada.

Otro hallazgo importante de este estudio es que las tareas domésticas se realizan de manera inconsistente debido a limitaciones de tiempo y energía, especialmente
entre aquellos con empleos de tiempo completo o estudiante con cargas académicas pesadas. Esta incosistencia es común en los hogares compartidos,
donde la carga se divide entre varias personas sin un sistema claro de asignación de responsabilidades o seguimiento. Sin herramientas adecuadas
para la organización, es posible que algunas tareas queden sin realizar o que recaigan siempre en las mismas personas, generando frustración y posibles 
conflictos.

El estudio también resalta que, a medida que los individuos se ven cada vez más presionados por sus estudios o trabajos, tienden a posponer las tareas del hogar, lo cual se 
agrava en entornos de convivencia compartida. Este aplazamiento puede llevar a una acumulación de tareas, que a su vez aumenta el nivel de estrés de los residentes y genera tensiones en el hogar.

Por último, los resultados también sugieren que la organización efectiva del hogar es crucial para el bienestar de sus miembros. En el contexto de hogares compartidos,
el uso de herramientas de planificación o aplicaciones específicas podría mejorar significativamente la administración de tareas y reducir el 
estrés relacionado con el hogar.

\newpage

\section{Estudio sobre los desafíos financieros en viviendas estudiantiles compartidas}
El estudio basado en datos de HUD, examina los desafíos financieros que enfrentan los estudiantes en viviendas compartidas. Uno de los hallazgos
principales de este estudio es que los contratos de arrendamiento de las viviendas a menudo requieren compromisos anuales. Los estudiantes deben 
seguir pagando el alquiler durante todo el año, incluso en períodos en que la vivienda está vacía, como los meses de verano. Este requisito genera una carga
financiera adicional, ya que los estudiantes deben encontrar formas de cubrir los pagos durante esos meses.

También señala que la organización de los gastos compartidos es una tarea compleja para los estudiantes. Dado que los contratos incluyen puede incluir 
diferentes servicios, la variación en las facturas mensuales añade dificultad al reparto equitativo de costes. Muchos estudiantes carecen de herramientas 
y sistemas de organización, lo que provoca confusión y conflictos en cuanto a las cantidades que cada miembro debe pagar. Esto puede ser 
díficil de gestionar cuando uno o varios de los residentes enfrentan problemas financieros, ya que retrasos o impagos afectan a todos los miembros del hogar.

Por último, este estudio también examina cómo las dificultades para gestionar los gastos afectan no solo la situación financiera de los estudiantes, 
sino también su bienestar general y rendimiento académico. La preocupación por la estabilidad financiera y la tensión asociada con la administración de los 
costos compartidos afectan a los estudiantes a nivel psicológico y académico, ya que deben dedicar su tiempo y energía a resolver estos problemas, lo que limita 
su concentración en los estudios.

\newpage

\section{Legislación actual de protección de datos}

La \textbf{Ley de Protección de Datos} está diseñada para garantizar el derecho a la privacidad y 
la protección de los datos personales, evitando el uso indebido o la divulgación de la información 
personal de personas. En la Unión Europea, la normativa más relevante es el 
\textbf{Reglamento General de Protección de Datos (GDPR)}, que entró en vigor en 2018 y 
establece un estándar estricto para el manejo de datos en todos los países miembros, aplicando también 
a empresas fuera de la UE que traten datos de ciudadanos europeos. En España, esta ley se complementa con 
la \textbf{Ley Orgánica de Protección de Datos y Garantía de los Derechos Digitales (LOPD-GDD)}, que adapta
el GDPR pero al contexto español.

Las principales medidas de este reglamento son:
\begin{enumerate}
    \item \textbf{Consentimiento explícito:} Las organizaciones deben obtener el consentimiento explícito de los usuarios antes de recopilar sus datos. Este consentimiento debe ser claro, específico e informado, lo que significa que los usuarios deben saber exactamente qué datos se recogen y para qué fines se usarán.
    \item \textbf{Derecho al Acceso, Rectificación y Supresión:} Los usuarios tienen el derecho de acceder a sus datos personales, solicitar correciones o eliminarlos. Esto se conoce como \textit{derecho al olvido} y permite a los usuarios mantener control sobre su información en línea.
    \item \textbf{Transparencia y Finalidad:} Los empresas deben informar a los usuarios sobre qué datos recopilan, el propósito de dicha recopilación y durante cuánto tiempo los conservarán. Esto busca evitar el almacenamiento de excesivo de datos y el uso para fines distintos a los incialmente indicados.
    \item \textbf{Seguridad de los datos:} Las organizaciones deben implementar medidas de seguridad para proteger los datos, como el cifrado y las políticas de acceso restringido. Estas medidas están orientadas a minimizar el riesgo de filtraciones de datos y accesos no autorizados.
    \item \textbf{Notificación de brechas de seguridad:} Si ocurre una violación de seguridad que afecte datos personales, las empresas deben notificar a las autoridades competentes y a los afectados en un plazo de 72 horas, asegurando así una respuesta rápida que mitigue posibles daños.
\end{enumerate}

Estas medidas buscan proteger los derechos de privacidad de los usuarios y obligan a las empresas a adoptar prácticas de manejo de datos que respeten la seguridad, confidencialidad y transparencia.





\end{document}